\section{Finite transfer and normalized limits}\label{sec:transfer}

We now transfer finitely many compatible coordinates to a smooth surface
with finitely many holes.  The transfer preserves the coordinates and the
original conformal structure near a prescribed compact core.  Outside the
core we choose the coordinates and the conformal structure together, so
that the added energy and the excess over the given oscillation budgets
are small.  The key is to make pairwise wedge products small: a suitable
metric then makes the sum
of the coordinate energies small as well.

After proving this finite construction, we normalize its circle-domain
maps and pass to limits on $G$.  Nested source collars will identify all
complementary components of the limit and locate the round disks
associated with prescribed dust points.

\subsection{The transfer theorem}

The construction selects finitely many disjoint dust disks.  Some are
filled, with a conformal structure chosen below; the others have their
interiors excised and are called \emph{terminal disks}.  Each terminal
disk is indexed by a distinct dust point in its interior.

In expressions involving a physical completion, we use the same symbol
$M_i$ for the inverse image $\pi^{-1}(M_i)$ of a mask in $X$.
Thus $Q\cap M_i$ below is a subset of the physical sphere.

\begin{theorem}[Finite transfer]\label{thm:transfer}
Let $u_1,\ldots,u_m$ be a finite ordered list satisfying
Definition~\ref{def:coordinates} and compatibility~(P), with a common
countable flag set $\Sigma$.  The same assertion holds for a finite
ordered sublist of a countable compatible list.  Prescribe a finite set
$F\subset B$ containing all singular points of the list, a compact core
$K_0\subset G$, a dust-disk cover from the fixed exhaustion, and
$\eps>0$.

There exist a smooth compact core $K\subset G$, finitely many
terminal disks $T_s$ with pairwise disjoint closures, indexed by a finite set
$S\subset B$ containing $F$, and a smooth metric $g$ on the physical
sphere, with the following properties.  Every terminal disk lies in a
member of the prescribed cover, $s\in T_s$, and $K$ contains $K_0$ and
all mask boundaries in its interior.  The metric $g$ agrees with the
spherical metric on a neighborhood of $K$.  Write
\[
 Q=\Chat\setminus\bigcup_{s\in S}T_s,
 \qquad Q^\circ=\Chat\setminus\bigcup_{s\in S}\overline{T_s}.
\]
Then $K\subset Q^\circ$, and $Q$, equipped with $g$, is a smooth
bordered surface of genus zero.
Each $u_i$ has a transferred coordinate $u_i^*$ on $Q\cap M_i$, smooth
up to its boundary in local real branches, with the following properties.
\begin{enumerate}[label=\textup{(\roman*)}]
\item $u_i^*=u_i$ on a neighborhood of $K\cap M_i$;
\item The total energy outside $K$ satisfies
\begin{equation}\label{eq:transfer-tail}
 \sum_{i=1}^m
 \|\dd u_i^*\|_{2,(Q\setminus K)\cap M_i,g}^2<\eps^2;
\end{equation}
\item For each nonsingular terminal index $s\in S\cap M_i$, a
single-valued lift on a collar of $\partial T_s$ has boundary oscillation
$A_i(T_s)$ satisfying
\begin{equation}\label{eq:transfer-oscillation}
 \left(\sum_{\substack{s\in S\cap M_i\\s\ne p_i}}
       [A_i(T_s)-a_i(s)]_+^2\right)^{1/2}<\eps.
\end{equation}
Here $[x]_+=\max\{x,0\}$.  For a real coordinate the restriction
$s\ne p_i$ is omitted.
\item A periodic coordinate still adds $1$ under its positive deck
translation.  Its differential has period $1$ around the terminal disk
indexed by its singular point, and period $0$ around every other
terminal disk in its mask.  Here a positive circuit around a terminal
disk is counterclockwise around that disk.
\end{enumerate}
The core $K_0$ may contain any prescribed compact subsets of $G$.
In particular the conformal structure on $\{|z|\ge R\}$ can always be
retained.  Terminal disks can be required to lie in arbitrarily small
prescribed dust neighborhoods of their indices.
\end{theorem}

\begin{proof}
We first replace the coordinates outside a large core by shifted
staircases.  Compatibility will provide disks on whose collars at most
one staircase varies.  We can then fill those disks so that the extended
differentials have rank at most one, and choose a metric from these
differentials.

We regard $\dd u_i$ as zero wherever its mask is absent; every use of
this convention is away from mask boundaries.  Fix $\eta>0$, to be
chosen in terms of $\eps$ at the end.

\smallskip\noindent
\emph{Step 1: random staircases.}
Choose a smooth core $K$ containing $K_0$, all mask boundaries, and the
exterior $\{|z|\ge R\}$ in its interior.  Enlarge it until the original
energies outside $K$ are so small that
\begin{equation}\label{eq:transfer-original-wedges}
 \sum_{i<j}\int_{G\setminus K}
       |\dd u_i\wedge\dd u_j|
 \le \sum_{i<j}
       \|\dd u_i\|_{2,G\setminus K}
       \|\dd u_j\|_{2,G\setminus K}
 <\frac{\eta}{8}.
\end{equation}
The mask intersections are implicit in these integrals.  Take
$\chi\in C^\infty(G,[0,1])$ equal to $0$ on a
neighborhood of $K$ and equal to $1$ outside a larger compact set.

The staircases will vary only near a finite grid of levels modulo one.
Let $h=1/N$, where $N$ is a positive integer.  Independently for each
$i$, choose $\theta_i$ uniformly in $[0,h)$.  For $0<\delta<h/4$ take a
smooth nondecreasing staircase $P_{i,\delta}$ with
\[
 P_{i,\delta}(x+h)=P_{i,\delta}(x)+h,
 \qquad |P_{i,\delta}(x)-x|\le h,
\]
whose derivative is supported within distance $\delta$ of
$\theta_i+h\Z$.  It is obtained by translating one fixed staircase
with an $h$-periodic derivative of integral $h$ over a period.  Therefore
\begin{equation}\label{eq:transfer-staircase-mean}
 \E P_{i,\delta}'(x)=1
 \quad\text{for every }x.
\end{equation}
Because $1=Nh$, it also satisfies
$P_{i,\delta}(x+1)=P_{i,\delta}(x)+1$.

Use the cutoff to join each staircase to the original coordinate:
\begin{equation}\label{eq:transfer-staircase}
 v_{i,\delta}
 =u_i+\chi\bigl(P_{i,\delta}(u_i)-u_i\bigr).
\end{equation}
The formula is interpreted in local branches and respects integer
translations, so it is consistent for periodic coordinates.  Set
\[
 q_{i,\delta}=1-\chi+\chi P_{i,\delta}'(u_i),
 \qquad r_{i,\delta}=P_{i,\delta}(u_i)-u_i.
\]
Then $q_{i,\delta}\ge0$, $|r_{i,\delta}|\le h$, and
\[
 \dd v_{i,\delta}
 =q_{i,\delta}\,\dd u_i+r_{i,\delta}\,\dd\chi.
\]
Equation~\eqref{eq:transfer-staircase-mean} and independence for distinct
labels give $\E q_{i,\delta}=1$ and
$\E(q_{i,\delta}q_{j,\delta})=1$ for $i\ne j$.  Consequently the wedge
sum
\[
 W_\delta
 =\sum_{i<j}\int_{G\setminus K}
             |\dd v_{i,\delta}\wedge\dd v_{j,\delta}|
\]
satisfies
\begin{equation}\label{eq:transfer-wedge-mean}
 \E W_\delta
 \le \sum_{i<j}\int_{G\setminus K}|\dd u_i\wedge\dd u_j|
   +(m-1)h\sum_i\int_{\supp\dd\chi}
                         |\dd\chi|\,|\dd u_i|\,\dd a.
\end{equation}
The last integrals are finite because $\supp\dd\chi$ is compact in
$G$.  With $K$ and $\chi$ fixed, choose $h$ small enough that the
right side is less than $\eta/4$, uniformly in $\delta$.

\smallskip\noindent
\emph{Step 2: choosing the grids and terminal indices.}
For the random shifts above, put
\[
 F_i^{\rm grid}
   =\bigcup_{k=0}^{N-1}H_i(\theta_i+kh).
\]
Each grid is compact.  In addition to $F$, retain as terminal indices
all flagged points lying in the grids of two distinct labels:
\[
 S=F\cup
 \bigcup_{i<j}\bigl(\Sigma\cap F_i^{\rm grid}
                                  \cap F_j^{\rm grid}\bigr),
 \qquad J=|S|.
\]
At a nonsingular dust point, a band of length at most $a_i(s)$ meets
the translated lattice with probability at most
$\min(1,a_i(s)/h)$.  The same bound holds at the singular point because
there $a_i(s)=1$.  Thus, by independence,
\begin{equation}\label{eq:transfer-terminal-mean}
 h^2\E J
 \le h^2|F|+
       \sum_{i<j}\sum_{s\in\Sigma}
                 \min(h,a_i(s))\min(h,a_j(s))
 \longrightarrow0\qquad(h\downarrow0).
\end{equation}
Each summand tends to zero and is bounded by $a_i(s)a_j(s)$, whose
sum is finite by Cauchy--Schwarz.  Dominated convergence proves the
limit.  The same estimate gives $J<\infty$ almost surely for each $h$.

Decrease $h$, if necessary, so that $h^2\E J<\eta^2/4$ as well.
Fix a deterministic sequence $\delta_k\downarrow0$ with
$\delta_k<h/4$.  Fatou's Lemma and
Equation~\eqref{eq:transfer-wedge-mean} give
\[
 \E\bigl(\liminf_{k\to\infty}W_{\delta_k}\bigr)<\eta/4.
\]
Markov's inequality shows that a set of positive probability satisfies
both
\begin{equation}\label{eq:transfer-grid-selection}
 \liminf_{k\to\infty}W_{\delta_k}<\eta,
 \qquad h\sqrt J<\eta.
\end{equation}

We also require compatibility~(P) for every choice of one lattice
level for each of any collection of distinct labels.  Each such tuple
has an absolutely continuous distribution in its level variables, and
there are only finitely many choices.  Hence this additional requirement
has probability one.  Fix grids satisfying it and
Equation~\eqref{eq:transfer-grid-selection}.

The grids and their finite set $S$ are now fixed.  Choose terminal
disks about the points of $S$, with pairwise disjoint closures, in the
region where $\chi=1$.  They contain no other
point of $S$, lie in the prescribed dust neighborhoods, and avoid all
mask boundaries.  For every nonsingular coordinate present at $s$, its
values in a sufficiently small neighborhood of $s$, in an appropriate
local lift, lie in its band enlarged by $h$ at each end.  We therefore
choose $T_s$ so that
\[
 \osc_{\partial T_s}u_i\le a_i(s)+2h.
\]
The bound $|P_{i,\delta}(x)-x|\le h$ gives, simultaneously for every
$\delta$,
\begin{equation}\label{eq:transfer-terminal-oscillation}
 \osc_{\partial T_s}v_{i,\delta}\le a_i(s)+4h.
\end{equation}
No real oscillation bound is required at a singular index.  The
terminal disks are now fixed, while the staircase width can still
decrease.

\smallskip\noindent
\emph{Step 3: filling disks whose boundaries meet at most one grid.}
We claim that the remaining dust has a finite cover by disks, with
pairwise disjoint closures and disjoint from the terminal disks, such
that the boundary of each new disk meets at most one
$F_i^{\rm grid}$.  All these disks can be placed where $\chi=1$,
inside the prescribed dust neighborhoods, and away from mask
boundaries.

First observe that if $s\in B\cap F_j^{\rm grid}$ is not a terminal
index, then $s$ is a singleton component of
\begin{equation}\label{eq:transfer-lower-obstacle}
 B\cup\bigcup_{i<j}F_i^{\rm grid}.
\end{equation}
If $s$ is flagged, it does not belong to any lower grid, by the choice
of $S$.  If $s$ is free, each coordinate defined at $s$ has a unique
trace.  Thus, near $s$, only one level from each of the grids occurring
at $s$ can be present; the other levels are closed and miss $s$.
If a nondegenerate component of the lower grid union contained $s$,
intersect that component with a sufficiently small closed neighborhood
of $s$.  Boundary bumping makes the component of this intersection
through $s$ nondegenerate.  It uses only the lower levels present at
$s$, contradicting the chosen instance of compatibility~(P).
Lemma~\ref{lem:dust-disks}\textup{(i)} shows that adjoining $B$ cannot
create another nondegenerate continuum through $s$.  This proves the
singleton assertion in both cases.

Process the labels in decreasing order.  At the stage for $j$, the
remaining dust is a compact relatively clopen subset of $B$, disjoint
from the already chosen closed disks and from every higher grid.  Use
Lemma~\ref{lem:dust-disks} to cover it by finitely many disjoint small
preliminary disks whose closures avoid these obstacles.  In each such
disk $V$, the remaining points of $F_j^{\rm grid}\cap B$ form a
compact set.  Each is a singleton component of the closed set in
Equation~\eqref{eq:transfer-lower-obstacle}; hence
Lemma~\ref{lem:dust-disks} encloses it in a smaller disk with closure
in $V$ and boundary avoiding that closed set.

Take finitely many such disks covering the compact target.  To merge
them while preserving boundary avoidance, perturb their
finitely many smooth boundaries into general position within the open
complement of the obstacle, take the outer contours of their union,
and smooth the corners there.  Fill bounded complementary regions and
discard nested contours.  Every resulting disk stays in $V$, because
the complement of $V$ is connected and contains infinity.  The
resulting closures are pairwise disjoint, their interiors still cover
the target, and their boundaries avoid all lower grids.  They also
avoid the higher grids, since their closures lie in $V$.  Add these
disks to the chosen collection and continue.  After the last label,
cover the dust still remaining by disks avoiding all the grids.  This
proves the claim.  Call all disks chosen in this step
\emph{filling disks}.

\smallskip\noindent
\emph{Step 4: smooth fillings of rank at most one.}
The finite collection of terminal and filling disks covers $B$ in its
interiors.  Its complement is therefore a compact subset of $G$.
On each filling boundary, every coordinate except possibly one avoids
its grid.  Choose connected collars, compact in $G$, on which this
avoidance still holds and $\chi=1$.  Every singular point lies in its
assigned terminal disk, so the coordinates have single-valued lifts on
each filling collar.  When the collar avoids $F_i^{\rm grid}$, a real
lift has connected image in one interval complementary to the level
lattice $\theta_i+h\Z$.  Compactness of the collar
then makes its staircase constant for all sufficiently small $\delta$.
Call these constant coordinates inactive on the collar.  There are
only finitely many collars, so one width bound works for all of them.

By Equation~\eqref{eq:transfer-grid-selection} we may choose one of the
$\delta_k$ small enough for all these collar conditions and with
$W_{\delta_k}<\eta$.  Fix this width and suppress it from the notation.
Retain $v_i$ on the complement of the selected disks.  Across each
filling disk extend the inactive coordinates constantly.  Extend the
one possibly varying coordinate from an inside collar by a smooth
cutoff, making it constant on a smaller interior disk.  These
extensions agree with the retained coordinates on full collars and
are therefore smooth.  Their differentials have rank at most one
at every point of a filling disk, so
\begin{equation}\label{eq:transfer-filling-wedges}
 \dd u_i^*\wedge\dd u_j^*=0
 \quad\text{on every filling disk.}
\end{equation}
Here $u_i^*$ denotes the resulting coordinate on the completed
bordered surface $Q$.

To check the periods, pull back by $\pi$ the cyclic cover of the
punctured mask.  This cover is trivial on each filling disk.  Perform
the extension on one sheet and use its integer translates on the
others.  Equation~\eqref{eq:transfer-staircase} respects the same
translations, so the completed coordinate has its original additive
period.  Under $\pi$, a counterclockwise circuit about $T_{p_i}$ winds
once about $p_i$, while a circuit about any other terminal disk in the
mask has winding number zero.  Thus the periods are exactly those
asserted in the theorem.

\smallskip\noindent
\emph{Step 5: choosing a metric.}
The coordinates are now extended, with small wedge integrals outside
$K$ and zero wedge products in the filling disks.  It remains to turn
this control into small energy in one smooth metric.
Let $g_0$ be the spherical metric and, outside $K$, form the smooth
positive semidefinite tensor
\[
 A=\sum_{i:\,M_i\text{ present}}\dd u_i^*\otimes\dd u_i^*.
\]
The sum at a point includes the masks containing that point.  All mask
boundaries lie inside $K$, so this tensor is smooth where it is used.
The differentials are single-valued even for periodic coordinates.
For $\tau>0$ put
$g_\tau=A+\tau g_0$.  In a $g_0$-orthonormal frame, let
$\lambda_1,\lambda_2\ge0$ be the eigenvalues of $A$.  The sum of the
energy densities of the coordinates in $g_\tau$, expressed relative
to $\dd a_{g_0}$, is
\begin{equation}\label{eq:transfer-metric-density}
 \operatorname{tr}(g_\tau^{-1}A)\sqrt{\det g_\tau}
 =\frac{2\lambda_1\lambda_2+\tau(\lambda_1+\lambda_2)}
        {\sqrt{(\lambda_1+\tau)(\lambda_2+\tau)}}
 \longrightarrow 2\sqrt{\lambda_1\lambda_2}.
\end{equation}
In particular the limit is zero when $A$ has rank at most one.
The Cauchy--Binet identity gives
\[
 \det A=\sum_{i<j}|\dd u_i^*\wedge\dd u_j^*|_{g_0}^2.
\]
The transferred differentials are bounded on the finitely many smooth
compact pieces of the completed surface.  The densities in
Equation~\eqref{eq:transfer-metric-density} therefore have a common
integrable bound for $0<\tau\le1$: each term
$\lambda_1\sqrt{(\lambda_2+\tau)/(\lambda_1+\tau)}$ is at most
$\sqrt{\lambda_1(\lambda_2+1)}$, and similarly for the other term.
Dominated convergence is therefore applicable.  By
Equation~\eqref{eq:transfer-filling-wedges} and the choice of the width,
\begin{equation}\label{eq:transfer-optimized-energy}
 \lim_{\tau\downarrow0}
       \sum_i\|\dd u_i^*\|_{2,(Q\setminus K)\cap M_i,g_\tau}^2
 =2\int_{Q\setminus K}\sqrt{\det A}\,\dd a_{g_0}
 \le2W_\delta<2\eta.
\end{equation}

Choose $\tau$ so that the total energy outside $K$ in $g_\tau$ is
less than $3\eta$.  With this regularization fixed, interpolate between
$g_0$ and $g_\tau$ in a thin exterior collar of $\partial K$, keeping
$g_0$ on a neighborhood of $K$.  The interpolated metrics are uniformly
positive definite and their energy densities are bounded on the
collar.  We may therefore choose its thickness so that the
interpolation adds less than $\eta$ to the total energy.  Extend the
resulting metric smoothly across each terminal cap, first across its
collar and then by interpolation with $g_0$ inside the cap.  This
does not change the metric on $Q$.

The tail energy is now less than $4\eta$.  Moreover,
Equation~\eqref{eq:transfer-terminal-oscillation} gives for each $i$
\[
 \left(\sum_{s\ne p_i}[A_i(T_s)-a_i(s)]_+^2\right)^{1/2}
 \le4h\sqrt J<4\eta.
\]
Taking $\eta<\min(\eps^2/4,\eps/4)$ establishes
Equations~\eqref{eq:transfer-tail} and
\eqref{eq:transfer-oscillation}.  All alterations occurred outside a
neighborhood of $K$, completing the proof.
\end{proof}

We will also need energy control away from a prescribed part of the
core.  Exact agreement on $K$ and the tail estimate give, for every
measurable $E\subset K\cap M_i$,
\begin{equation}\label{eq:transfer-localized-energy}
 \|\dd u_i^*\|_{2,(Q\cap M_i)\setminus E,g}^2
 \le \|\dd u_i\|_{2,(G\cap M_i)\setminus E}^2+\eps^2.
\end{equation}
Both this estimate and the boundary oscillation estimates are unchanged
by a conformal uniformization of $Q^\circ$.

\subsection{Circle-domain completions and compactness}

We now apply finite uniformization to these completions and obtain a
subsequential limit on the original domain.  Keeping a fixed exterior
region unchanged will give one normalization and one bound for all
the finite maps.

Extend the metric of Theorem~\ref{thm:transfer} across the terminal
caps, as in its proof.  Compact genus-zero uniformization identifies
this metric sphere conformally with $\Chat$
\cite[Theorem~10-3]{Ahlfors1973}.  Classical finite
circle-domain uniformization then maps $Q^\circ$ conformally onto a
circle domain; only the finite-connectivity theorem is used here
\cite{Koebe1920}; see also \cite[Theorem~15.7.9]{Conway1995}.
Postcompose by a M\"obius transformation sending the image of infinity
to infinity, and then by an affine map, to obtain
\begin{equation}\label{eq:transfer-normalization}
 f_Q(z)=z+O(1/z).
\end{equation}
These postcompositions preserve roundness.

In the target coordinate $w$ we use Euclidean length $|\dd w|$ and
Euclidean area; on the source we continue to use physical spherical
measures.  Conformal invariance makes the Dirichlet energy norms agree.

Fix a countable compatible list and a countable set of dust points,
called \emph{marks}, containing all its singularities.  Choose an enumeration of the underlying countable set of labels.
The enumeration is used only to choose increasing finite subsets;
each subset is transferred in the order inherited from the original
well-order.  It need not be an initial segment of that well-order.
This distinction allows a list of any countable ordinal type, including
orders greater than $\omega$.  Apply
Theorem~\ref{thm:transfer} to these finite ordered sublists and
increasing finite sets of marks.  Prescribe cores containing an exhaustion of
$G$, let the errors tend to zero, and place all terminal disks in
members of dust covers whose mesh tends to zero.  Denote the normalized
map of the $n$th completion by $f_n$.

Each $f_n$ is defined on its completed physical surface.  The conformal
structure there can differ from the original structure outside the
retained core.  For every fixed compact subset of $G$, the maps are
eventually holomorphic and injective on a neighborhood of that subset
in the original structure.  Each fixed coordinate is eventually
transferred, agrees with its original values on exhausting cores,
and has errors tending to zero in
Equations~\eqref{eq:transfer-tail},
\eqref{eq:transfer-oscillation}, and
\eqref{eq:transfer-localized-energy}.

Every fixed mark $b$ eventually indexes a terminal disk shrinking to
$b$ in $X$.  Write $D_{n,b}$ for its round complementary disk in the
target.  A finite completion may also have terminal indices outside
the fixed countable set of marks; these additional indices may depend
on $n$.

\begin{lemma}[Normalized compactness]\label{lem:normalized-limits}
There is a constant $L$, depending only on $R$, with the following
properties.  For every completion just described, all complementary
disks of its normalized circle domain and all images of physical
points in its domain with $|z|\le4R$ lie in $\{|w|<L\}$.
One may take $L=5R$.
After passage to a subsequence, $f_n$ converges locally uniformly on
$G$ to a normalized univalent map $f$.  Simultaneously, for every
fixed mark $b$, the disks $D_{n,b}$ converge in the Hausdorff metric to
a closed round disk $D_b$, possibly a singleton.

The same bound $L$ holds for the bounded complementary components and
the images of $G\cap\{|z|\le4R\}$ for every normalized univalent map
on $G$.
\end{lemma}

\begin{proof}
The exterior $\{|z|\ge R\}$ retains its original conformal
structure.  Write
\[
 \frac{f_n(R\zeta)}R
 =\zeta+\sum_{j\ge1}c_{n,j}\zeta^{-j},\qquad |\zeta|>1.
\]
The exterior univalent area theorem gives
$\sum_{j\ge1}j|c_{n,j}|^2\le1$
\cite[Theorem~5-1]{Ahlfors1973}.  Hence, for $|\zeta|=r>1$,
\begin{equation}\label{eq:transfer-exterior-bound}
 \left|\frac{f_n(R\zeta)}R-\zeta\right|
 \le\left(\sum_{j\ge1}\frac{r^{-2j}}j\right)^{1/2}
 =\sqrt{-\log(1-r^{-2})}.
\end{equation}
For $r\ge2$ the right side is at most
$\sqrt{\log(4/3)}<1$.

The image of $\{|z|>2R\}\cup\{\infty\}$ is the unbounded side of
the Jordan curve $f_n(\{|z|=2R\})$.  Injectivity therefore puts all
other image points and all complementary disks on its bounded side,
which lies in $\{|w|<3R\}$ by
Equation~\eqref{eq:transfer-exterior-bound}.  For
$2R\le|z|\le4R$ that equation gives $|f_n(z)|<5R$.  This proves the
uniform bound.  The same argument applies without change to any
normalized univalent map on $G$.

On every relatively compact connected subdomain of
$G\setminus\{\infty\}$, the maps are eventually holomorphic,
injective, and uniformly bounded.  Diagonal normal-family compactness
therefore gives a locally uniform holomorphic limit.  The exterior
estimate preserves its Laurent normalization, so the limit extends
across infinity as $f(z)=z+O(1/z)$ and is nonconstant.  Hurwitz's
theorem for limits of injective holomorphic maps, applied on each of
these subdomains, makes the limit injective on all of $G$.

The centers and radii of the marked disks lie in a fixed compact
parameter set.  A further diagonal subsequence makes them converge
for all countably many marks.  Convergence of centers and nonnegative
radii is precisely Hausdorff convergence of the disks, with radius
zero giving a singleton.
\end{proof}

\subsection{Identification and trapping of the ends}

We have obtained a univalent limit and round disks at the marked
ends.  To compare those disks with the full complement of the image,
we identify each complementary component by nested source collars.

For a normalized univalent map $f$ on $G$ and a smooth Jordan curve
$C\subset G$ enclosing a dust disk, write
$\mathcal Q_f(C)$ for the filled closed bounded Jordan disk with
boundary $f(C)$.

\begin{lemma}[End components]\label{lem:end-components}
Let $f$ be a normalized univalent map on $G$.  For $b\in B$ set
\begin{equation}\label{eq:transfer-end-component}
 K_b(f)=\bigcap_{j\ge1}\mathcal Q_f(C_j(b)).
\end{equation}
These sets are exactly the complementary components of $f(G)$, with
one component for each $b\in B$.  They are nonempty continua, and
\[
 K_b(f)\subset\operatorname{int}\mathcal Q_f(C_j(b))
 \quad\text{for every }j.
\]
For a sequence and marked disk limits as in
Lemma~\ref{lem:normalized-limits},
\begin{equation}\label{eq:transfer-mark-trapping}
 D_b\subset K_b(f).
\end{equation}

In a single finite completion, each terminal boundary corresponds to
exactly one complementary round disk.  Points of the circle domain
approaching that disk have preimages approaching the corresponding
terminal boundary as a set.  Thus the cluster oscillation of a transferred
coordinate at a nonsingular disk is at most its boundary oscillation
$A_i(T_s)$, in a single-valued lift on a collar.
\end{lemma}

\begin{proof}
We first identify the components for an arbitrary normalized univalent
$f$, then treat terminal boundaries in a finite completion and pass
their marked disks to the limit.

Let $K_j$ be a core from the fixed exhaustion.  It is the closed
exterior of finitely many disjoint dust disks and contains infinity
in its interior.  Its image under $f$ is a compact embedded bordered
surface whose boundary curves are the images of the source boundary
curves.  Orientation fixes the image side of each boundary in a
collar.  The connected interior of the core, which contains infinity,
must remain on that exterior side.  Jordan separation therefore
identifies $f(K_j)$ with the closed exterior of the corresponding
pairwise disjoint filled bounded Jordan disks.

Under refinement of the cores these image disks are nested in exactly
the same way as their source dust disks.  Strict nesting in the source
gives
\[
 \mathcal Q_f(C_{j+1}(b))
 \subset\operatorname{int}\mathcal Q_f(C_j(b)).
\]
Their intersection is a nonempty compact continuum, strictly inside
each earlier image disk.  It misses $f(G)$, because every point of
$G$ lies in the interior of some exhaustion core, while the
intersection lies outside all image cores.

Conversely, a point outside $f(G)$ lies inside exactly one image dust
disk at every exhaustion level.  Refinement makes the corresponding
source disks nested.  Their mesh tends to zero, so they determine a
unique point $b\in B$, and the target point belongs to $K_b(f)$.
A connected subset of the complement cannot meet two disjoint image
disks at the same level.  Thus the component containing a point of
$K_b(f)$ lies in every disk in
Equation~\eqref{eq:transfer-end-component}.  The intersection is
itself connected, so it is the whole component.  Distinct dust points
are separated at some exhaustion level.  This proves the bijection.

For a finite completed surface, use cores bounded by successively
smaller collars of its terminal boundaries.  The same argument gives
one end per terminal boundary; its nested image disks intersect in
the corresponding round complementary disk.  Fix such a disk $D$
and a source collar curve $C$.  The disk $D$ is compactly contained
in the bounded side of the image of $C$.  Hence target points whose
distance from $D$ tends to zero eventually lie on that side, and their
preimages lie between $C$ and the terminal boundary.  Taking
arbitrarily small collars shows that the preimages approach this
boundary as a set.

The transferred coordinate is smooth up to the terminal boundary, so
its cluster values lie in its boundary range.  At a nonsingular
terminal index its period is zero, and a single-valued lift can be
used on the whole collar.  This gives the cluster oscillation bound
without a boundary extension of the uniformizing map.

Finally fix a mark $b$ and an exhaustion curve $C_j(b)$.  For all
sufficiently large $n$ this curve is retained on the core and the
terminal disk indexed by $b$ lies on its bounded source side.  Jordan
separation in that finite completed surface gives
\[
 D_{n,b}\subset\operatorname{int}\mathcal Q_{f_n}(C_j(b)).
\]
Uniform convergence on the curve and stability of winding number put
every Hausdorff limit point of these disks in
$\mathcal Q_f(C_j(b))$.  Intersecting over $j$ proves
Equation~\eqref{eq:transfer-mark-trapping}.  For strict trapping inside
the image of a fixed collar, first use $C_{j+1}(b)$ and then its
compact containment in the bounded side of $f(C_j(b))$.
\end{proof}

The later selections will enlarge the coordinate list, flags, and
marks.  Every countable compatible enlargement still admits the
construction above.  Each fixed coordinate and each fixed mark are
eventually included in its diagonal transfers, so all their estimates
and the trapping conclusions persist.
